\section{The original model and calculus}
\label{sec:original}

This section recalls Chapter~4 of the dissertation, in its notation. Nothing
here is new. Page and figure numbers refer to the dissertation.

\subsection{The machine}

A \emph{memory system} is a triple $\mu = (h,B,K)$: a heap
$h : \setlocations \pfun \setvalues$, an array of write buffers
$B : \setprocessors \tfun \listof{(\setlocations \times \setvalues)}$ with
the oldest write first, and a set $K \subseteq \setprocessors$ of
\emph{blocked} processors. A state pairs a stack $s$ with a memory system.
The lock is not a separate component: it is free when $K = \nil$ and held
by $i$ when $K = \setprocessors \setminus \set{i}$.

Primitive commands carry the identifier of the processor that executes
them. Their behaviour on processor $i$ is as follows.
\begin{itemize}[nosep]
  \item $\cstore{e}{e'}_i$ appends a write to $B(i)$. The address must be
        allocated, in the heap or in $B(i)$. A blocked processor may still
        store.
  \item $\cload{x}{e}_i$ requires $i \notin K$ and returns the newest write
        to that address in $B(i)$, or else the heap value.
  \item $\cfence_i$, $\clock_i$ and $\cunlock_i$ require $B(i) = \lnil$.
        $\clock_i$ additionally requires $K = \nil$ and sets
        $K = \setprocessors \setminus \set{i}$; $\cunlock_i$ reverses this.
  \item A silent step $\taustep$ moves the oldest write of some unblocked
        processor from its buffer to the heap.
\end{itemize}
Commands are built from primitives by sequencing $\cseq{c}{c'}$, choice
$\cchoice{c}{c'}$, iteration $\cloop{c}$ and parallel composition
$\cpar{c}{c'}$. There is no allocation. While a processor holds the lock,
no other processor can load or flush, but others may continue to store and
to assign to registers.

\subsection{Generalized memories and separation}

A \emph{generalized memory system} $\nu = (h,B,K,\Gamma)$ adds a
\emph{buffer-completeness} set $\Gamma \subseteq \setprocessors$, with
$h = \nil$ whenever $\Gamma = \nil$. Membership $i \in \Gamma$ means that
buffer $i$ is complete: nothing may be placed before it.

Three operations combine generalized memories. Write $\nu_k =
(h_k,B_k,K_k,\Gamma_k)$. All three take $K_0 \cup K_1$ and
$\Gamma_0 \cup \Gamma_1$ and require $K_0 \cap K_1 = \nil$.
\begin{description}[nosep,leftmargin=1.5em]
  \item[Spatial, $\nu_0 \msep \nu_1$.] The address domains of the two
    sides, counting both heaps and buffers, are disjoint. The heap is
    $h_0 \uplus h_1$ and each buffer is an interleaving of the two.
  \item[Temporal, $\nu_0 \mseq \nu_1$.] The heap is $h_0 \override h_1$, so
    later values win, and each buffer is $B_0(i) \lapp B_1(i)$. The sides
    may share addresses, subject to a compatibility condition: an unblocked
    pending write in $\nu_0$ may not target an address in the heap of
    $\nu_1$, nor an address pending on a different unblocked processor in
    $\nu_1$. Finally $B_0(i) = \lnil$ for every $i \in \Gamma_1$.
  \item[Spatiotemporal, $\nu_0 \mhash \nu_1$.] As temporal, but each buffer
    is any interleaving whose lookup function agrees with
    $B_0(i) \lapp B_1(i)$, and the condition on $\Gamma_1$ is dropped. It
    contains both of the others and is used only in the semantics of
    specifications.
\end{description}
\emph{Strong temporal} separation $\nu_0 \msseq \nu_1$ is $\nu_0 \mseq
\nu_1$ restricted to pairs that are also spatially compatible.

\subsection{Assertions}

Assertions are given by
\begin{align*}
  P \bnfdef{} & b \bnfbar P \disj P \bnfbar P \conj P \bnfbar \exists x \st P
    \bnfbar \forall x \st P \bnfbar \femp \bnfbar \fbar{e} \bnfbar
    \flock{e} \bnfbar e \fwrite{e'} e'' \\
  \bnfbar{} & P \fsep P \bnfbar P \fseq P \bnfbar P \fsepi P \bnfbar
    P \fseqli P \bnfbar P \fseqri P
\end{align*}
and are interpreted over a stack and a generalized memory. The atoms are:
\begin{itemize}[nosep]
  \item $\femp$: empty heap, empty buffers, $K = \nil$; $\Gamma$ is
        unconstrained.
  \item $\fbar{i}$: as $\femp$, with $i \in \Gamma$ when $i$ is a processor.
  \item $\flock{i}$: empty heap and buffers, with
        $K = \setprocessors \setminus \set{i}$.
  \item $e \fwrite{i} e'$ (``leads to''): either a single \emph{pending}
        write $(e,e')$ in $B(i)$ with an empty heap and
        $K \subseteq \set{i}$, or a single \emph{flushed} heap cell
        $e \mapsto e'$ with empty buffers, $K = \nil$ and $i \in \Gamma$.
\end{itemize}
$P \fsep Q$ and $P \fseq Q$ hold of $\nu$ when $\nu$ is a spatial,
respectively temporal, combination of a model of $P$ and a model of $Q$.
The three implications are the corresponding adjoints. Two abbreviations
matter below:
\[
  P \fsseq Q \eqdef (P \fsep Q) \conj (P \fseq Q)
  \qquad\qquad
  e \fpointsto e' \eqdef (\exists x \st e \fwrite{x} e') \fseq
                         (\forall x \st \fbar{x}).
\]
Points-to therefore describes a flushed cell with \emph{every} processor
buffer-complete. The dissertation claims that every assertion is closed
under flushing (Proposition~9), and that each conjunction is associative
with $\femp$ as a unit (Figure~4.5).

\subsection{Specifications and rules}

A specification $\spec{J}{P}{c}{Q}$ has an invariant $J$ for shared state
and a private precondition and postcondition. Figure~\ref{fig:original}
reproduces the axioms and rules.

\begin{figure}[t]
\small
\begin{mathpar}
\inferrule[skip]{ }{\spec{J}{P}{\cskip_i}{P}}
\and
\inferrule[assume]{ }{\spec{J}{\neg b \disj P}{\cassume{b}_i}{P}}
\and
\inferrule[assert]{ }{\spec{J}{b \conj P}{\cassert{b}_i}{P}}
\and
\inferrule[assign]{ }{\spec{J}{P\subst{e}{x}}{\cassign{x}{e}_i}{P}}
\and
\inferrule[fence]{ }{\spec{J}{\femp}{\cfence_i}{\fbar{i}}}
\and
\inferrule[lock]{ }{\spec{J}{\femp}{\clock_i}{\flock{i}}}
\and
\inferrule[unlock]{ }{\spec{J}{\flock{i}}{\cunlock_i}{\femp}}
\and
\inferrule[load]{ }
  {\spec{J}{e \fwrite{i} e' \fsseq P}{\cload{x}{e}_i}
        {(e \fwrite{i} e' \fsseq P) \conj x = e'}}
\and
\inferrule[store]{ }
  {\spec{J}{e \fwrite{i} e'' \fsseq P}{\cstore{e}{e'}_i}
        {(e \fwrite{i} e'' \fsseq P) \fseq e \fwrite{i} e'}}
\and
\inferrule[frame-$\bullet$]
  {\spec{J}{P}{c}{Q} \\ \mod{c} \cap \fv{R} = \nil}
  {\spec{J}{R \fgen P}{c}{R \fgen Q}}
\and
\inferrule[cons]
  {P \sentails P' \\ \spec{J}{P'}{c}{Q'} \\ Q' \sentails Q}
  {\spec{J}{P}{c}{Q}}
\and
\inferrule[seq]
  {\spec{J}{P}{c}{R} \\ \spec{J}{R}{c'}{Q}}
  {\spec{J}{P}{\cseq{c}{c'}}{Q}}
\and
\inferrule[choice]
  {\spec{J}{P}{c}{Q} \\ \spec{J}{P}{c'}{Q}}
  {\spec{J}{P}{\cchoice{c}{c'}}{Q}}
\and
\inferrule[loop]{\spec{J}{P}{c}{P}}{\spec{J}{P}{\cloop{c}}{P}}
\and
\inferrule[disj]
  {\spec{J}{P}{c}{Q} \\ \spec{J}{P'}{c}{Q}}
  {\spec{J}{P \disj P'}{c}{Q}}
\and
\inferrule[ex]
  {\spec{J}{P}{c}{Q} \\ x \notin \fv{c,Q}}
  {\spec{J}{\exists x \st P}{c}{Q}}
\and
\inferrule[conc]
  {\spec{J}{P_k}{c_k}{Q_k} \\ \fv{P_k,c_k,Q_k} \cap \mod{c_{1-k}} = \nil
   \quad (k = 0,1)}
  {\spec{J}{P_0 \fsep P_1}{\cpar{c_0}{c_1}}{Q_0 \fsep Q_1}}
\and
\inferrule[share]
  {\spec{J}{P}{c}{Q}}
  {\spec{\femp}{J \fsep P}{c}{J \fsep Q}}
\and
\inferrule[inv]
  {\spec{\femp}{J \fsep P \fsep \flock{i}}{c}{J \fsep Q \fsep \flock{i}}}
  {\spec{J}{P \fsep \flock{i}}{c}{Q \fsep \flock{i}}}
\end{mathpar}
\caption{The original axioms and rules (Figures~4.7 and~4.8).
\rn{frame-$\bullet$} stands for three rules, \rn{frame-sp}, \rn{frame-tm}
and \rn{frame-stm}, with $\bullet$ read as $\fsep$, $\fseq$ and $\fsseq$.}
\label{fig:original}
\end{figure}

Three features of the design are worth recalling, because the repair turns
on them. First, \rn{lock} and \rn{unlock} produce and consume only the
token $\flock{i}$. The invariant reaches the lock holder through \rn{inv},
which may be applied to any command executed while the lock is held and
allows the private state to exchange cells with $J$. Second, \rn{frame-tm}
has only the usual variable condition: any history $R$ may be placed before
any specification. Third, temporal conjunction is not cancellative, as the
dissertation remarks: $\mu \sseq \mu_0 = \mu \sseq \mu_1$ does not imply
$\mu_0 = \mu_1$.

\subsection{Semantics and status}

Truth of a specification is defined through a family of predicates
$\safe{n}{c,s,\mu,J,Q}$ after Vafeiadis. The predicate holds at $n+1$ when:
if $c$ is $\cskip$ then $(s,\mu,\setprocessors) \sentails Q$; and for every
$\mu_J$, $\mu_F$ and $\mu_0 \in \mu_J \ssep (\mu_F \shash \mu)$ that is
lock-complete, with either $(s,\mu_J,\setprocessors) \sentails J$ or the
lock held in $\mu_0$, the configuration does not abort, and every step to
$(c',s',\mu_1)$ admits a new decomposition
$\mu_1 \in \mu'_J \ssep (\mu'_F \shash \mu')$ with
$\mu'_F \taurefines \mu_F$, the same condition on $J$, and
$\safe{n}{c',s',\mu',J,Q}$. A specification is true,
$\truespec{J}{P}{c}{Q}$, when $(s,\mu,\setprocessors) \sentails P$ implies
$\safe{n}{c,s,\mu,J,Q}$ for all $n$. Local states are always read at full
completeness, $\Gamma = \setprocessors$.

Theorem~1 states that provable specifications are true. The dissertation
gives soundness lemmas for the axioms in an appendix, sketches a frame lemma for the
spatiotemporal operator $\shash$, and reduces the spatial frame rule to Conjecture~1, which is left open.
The conjecture says that spatial separation survives a step: if a memory
splits spatially into $\mu_b$ and $\mu_c$ and takes a step, then any
spatiotemporal split of the result into a flushed descendant of $\mu_b$ and
some $\mu'_c$ is again a spatial split. Associativity of buffer overriding
(Proposition~2) and of the three separations, and flushing closure of
assertions (Proposition~9), are stated with at most a sketch.
